% ═══════════════════════════════════════════════════════════════════════════
%  پیش‌گفتار
% ═══════════════════════════════════════════════════════════════════════════
\chapter*{پیش‌گفتار}
\addcontentsline{toc}{chapter}{پیش‌گفتار}
\markboth{پیش‌گفتار}{پیش‌گفتار}

تاریخ بشر مملو از لحظاتی است که ملتی، قومی یا گروهی اجتماعی، ناتوان از رهایی خود به نیروی درونی، چشم به بیرون دوخته و از نیرویی خارجی یاری طلبیده است. این پدیده — که در این کتاب از آن با عنوان \concept{نجات‌طلبی} یاد می‌کنیم — پدیده‌ای چندلایه و پیچیده است که نمی‌توان آن را صرفاً در قالب «خیانت» یا «واقع‌گرایی» خلاصه کرد.

از \histref{شاهزاده بسوس}{۳۳۰ ق.م} که در برابر اسکندر مقدونی ایستاد تا \histref{شیخ خزعل}{۱۳۰۳ ش.} که با بریتانیا گره خورد؛ از \histref{ویلیام اورانژ}{۱۶۸۸ م.} که پارلمان انگلستان از او دعوت کرد تا \histref{مداخله‌ی ناتو در لیبی}{۲۰۱۱ م.} که شورای امنیت آن را تجویز کرد — همه‌ی این‌ها صورت‌بندی‌های مختلف یک پرسش بنیادین‌اند:

\begin{histQuote}{پرسش بنیادین}
آیا دعوت از نیروی خارجی برای رهایی از ستم داخلی، عملی مشروع و مؤثر است یا آغاز دوری تازه از وابستگی و خسران؟
\end{histQuote}

این کتاب تلاشی است برای پاسخ‌گویی علمی و مستند به این پرسش. رویکردم \concept{میان‌رشته‌ای} است: از تاریخ و علوم سیاسی گرفته تا حقوق بین‌الملل، جامعه‌شناسی، روان‌شناسی اجتماعی و اقتصاد سیاسی. تمرکز جغرافیایی ویژه بر \keyword{ایران} است — سرزمینی که در طول تاریخ بارها با این پدیده مواجه بوده — اما نگاهم جهانی و تطبیقی خواهد بود.

\section*{ساختار کتاب}

کتاب در چهار بخش اصلی و چهار پیوست سازمان یافته است:

\begin{itemize}[nosep, rightmargin=1em]
  \item \textbf{بخش نخست} (فصل‌های ۱ تا ۳): مبانی نظری، حقوقی و الگوهای تحلیلی.
  \item \textbf{بخش دوم} (فصل‌های ۴ تا ۷): بررسی تاریخی ایران از دوران باستان تا معاصر.
  \item \textbf{بخش سوم} (فصل‌های ۸ تا ۱۰): مداخلات خارجی در اروپا، جهان استعماری و دوران معاصر.
  \item \textbf{بخش چهارم} (فصل‌های ۱۱ تا ۱۴): تحلیل‌های روان‌شناسی اجتماعی، جامعه‌شناسی سیاسی، فعالیت‌های میسیونری-ایدئولوژیک و جمع‌بندی نهایی.
\end{itemize}

\section*{روش‌شناسی}

روش تحقیق \concept{توصیفی-تحلیلی} و \concept{تطبیقی-تاریخی} است. منابع شامل اسناد آرشیوی، پژوهش‌های آکادمیک معتبر به زبان‌های فارسی، انگلیسی، عربی و فرانسه، اسناد رسمی سازمان ملل، و اسناد آزادشده‌ی سازمان‌های اطلاعاتی است. ارجاع‌دهی در سراسر متن از نظام پانوشت استفاده می‌کند و کتاب‌نامه‌ی کامل در پایان کتاب آمده است.

\section*{نکته‌ی مهم درباره‌ی بی‌طرفی}

این اثر با رویکرد \concept{بی‌طرفانه‌ی علمی} نگاشته شده است. نویسنده تلاش کرده تمامی دیدگاه‌ها — از مدافعان مداخله‌ی انسان‌دوستانه تا منتقدان امپریالیسم — را منصفانه بازنمایی کند. با این حال، در پیوست «ب» یک نقد شخصی از منظر «گر رسم شود که مست گیرند، در شهر باید هر آنکه هست گیرند» ارائه شده است. همچنین نقدی به این نقد، از منظر «آیا بالاتر از سیاهی رنگی هست!» نگاشته شده است. 

\vspace{1cm}
\begin{flushleft}
  مهدی سالم\\
  تابستان ۱۴۰۴
\end{flushleft}

\cleardoublepage